%%═══════════════════════════════════════════════════════════════════
%% Lecture 01 — Introduction: Roundoff Errors, Numerical
%%              Differentiation, and Interpolation
%% 77315 — Computational Physics
%% Date: 26/03/2026
%%═══════════════════════════════════════════════════════════════════

\section{Lecture 01 — Introduction: Roundoff Errors, Numerical
  Differentiation, and Interpolation}\label{sec:lec01}

\begin{center}
\begin{tabular}{ll}
  \textbf{Date:}\quad 26/03/2026 & \\
\end{tabular}
\end{center}
\hrule\vspace{1.5em}

%%──────────────────────────────────────────────────────────────────
\subsection*{Overview}

Computational physics combines numerical analysis methods to solve physical
problems. Historically, physics progressed through a binary of theoretical
physics (creating mathematical models) and experimental physics (measuring
reality). Today, the computational approach serves as a formidable third
branch. By leveraging the immense processing power of computers, we can
solve complex models (such as nonlinear differential equations with no
analytical solutions) and perform macroscopic simulations of many-body
systems that would be otherwise impossible to track.

When translating pure mathematics into computer code, we must abandon the
notion of ``infinite precision.'' Three main factors must always be
considered:
\begin{enumerate}
    \item \textbf{Roundoff Errors}: Inaccuracies due to the finite memory
    computers use to store numbers.
    \item \textbf{Truncation Errors (Accuracy)}: Inaccuracies arising from
    using a finite mathematical approximation (like stopping a Taylor series
    early).
    \item \textbf{Stability}: The tendency of an algorithm to either dampen
    or exponentially amplify these microscopic errors over time.
\end{enumerate}

%%──────────────────────────────────────────────────────────────────
\subsection{Errors, Accuracy, and Stability}

\subsubsection{Catastrophic Cancellation and Roundoff Errors}
Computers represent real numbers using a limited number of bits. A common
representation is 4 bytes (Float32, yielding $\sim 7$ decimal digits of
precision) or 8 bytes (Float64, yielding $\sim 16$ decimal digits).

Because of this finite representation, every basic arithmetic operation
introduces a tiny relative error. While this is usually negligible, a severe
problem known as \textbf{catastrophic cancellation} occurs when subtracting
two very close numbers.

\begin{tcolorbox}[colback=blue!5!white,colframe=blue!75!black,
  title=Pedagogical Example: The Quadratic Formula]
Consider the quadratic equation $x^2 - 1000x + 0.001 = 0$. Here, $b = 500$
and $c = 0.001$. The standard quadratic formula gives the roots as
$x = b \pm \sqrt{b^2 - c}$.

If we calculate the smaller root using subtraction:
$x_{\text{small}} = 500 - \sqrt{250000 - 0.001} \approx 500 - 499.999999 =
0.000001$.

If our computer only has 7 digits of precision (Float32),
$\sqrt{250000 - 0.001}$ will be rounded to exactly $500.0000$. The
subtraction $500.0000 - 500.0000$ yields $0$. We have completely lost the
true answer due to catastrophic cancellation!

\textbf{The elegant solution}: Multiply the numerator and denominator by the
conjugate $b + \sqrt{b^2-c}$. This transforms the formula to
$x_{\text{small}} = \frac{c}{b + \sqrt{b^2 - c}}$.
Now: $x_{\text{small}} = \frac{0.001}{500 + 500} = \frac{0.001}{1000} =
10^{-6}$. By avoiding the subtraction of close numbers, we maintain perfect
precision.
\end{tcolorbox}

\subsubsection{Truncation Error and Stability}
Truncation error is the difference between the exact mathematical formula and
the approximated algorithm. For example, a Taylor series is infinite, but a
computer can only sum $N$ terms. The discarded terms constitute the truncation
error.

\textbf{Stability} dictates whether an algorithm is usable. An algorithm is
unstable if tiny roundoff or truncation errors grow exponentially as the
computation proceeds.

\begin{tcolorbox}[colback=red!5!white,colframe=red!75!black,
  title=Intuitive Analogy: The Golden Mean Recursion]
Consider calculating the powers of the ``Golden Mean'' $\phi =
\frac{\sqrt{5}-1}{2} \approx 0.618$. Mathematically, $\phi$ satisfies the
recurrence relation: $\phi^{n+1} = \phi^{n-1} - \phi^n$. This seems like a
brilliantly fast way to compute $\phi^{100}$ without multiplication!

However, this linear recurrence relation has two mathematical roots:
$\phi_1 \approx 0.618$ and $\phi_2 = \frac{-\sqrt{5}-1}{2} \approx -1.618$.
Even if we start with the exact values for $n=1, 2$, the tiny roundoff error
introduced by the computer acts as a ``seed'' for the second root. Because
$|\phi_2| > 1$, its contribution grows exponentially. By $n=50$, the
parasitic $\phi_2$ term completely dominates, and the algorithm blows up to
infinity. This is the hallmark of numerical instability.
\end{tcolorbox}

%%──────────────────────────────────────────────────────────────────
\subsection{Numerical Differentiation}
How do we teach a computer to calculate a derivative,
$\lim_{h \to 0} \frac{f(x+h)-f(x)}{h}$? We cannot set $h=0$, nor can we
make $h$ infinitely small due to roundoff errors.

If we have a discrete grid of points $x_n = x_0 + nh$, we can approximate
derivatives using Taylor expansions:
\begin{align*}
f(x_0 + h) &= f_0 + h f' + \frac{h^2}{2} f'' + \frac{h^3}{6} f''' +
\mathcal{O}(h^4) \\
f(x_0 - h) &= f_0 - h f' + \frac{h^2}{2} f'' - \frac{h^3}{6} f''' +
\mathcal{O}(h^4)
\end{align*}
Subtracting the second from the first cancels the $f''$ term, yielding the
\textbf{Central Difference Formula}:
\begin{equation}
f' = \frac{f_{+1} - f_{-1}}{2h} + \mathcal{O}(h^2)
\end{equation}
Geometrically, instead of drawing a secant line from $x$ to $x+h$ (which is
biased), we draw a secant line from $x-h$ to $x+h$, which perfectly balances
the slope around $x$.

\textbf{The Step-Size Dilemma.} If we make $h$ large, the Taylor
approximation fails (truncation error $e_t \sim h^2 f'''$). If we make $h$
microscopic, catastrophic cancellation destroys the numerator (roundoff error
$e_r \sim \frac{\eps_f f}{h}$). Minimising $E(h) = \frac{\eps_f f}{h} +
h^2 f'''$ yields:
\begin{equation}
h_{\text{opt}} \sim \sqrt[3]{\frac{3\,\eps_f\, f}{f'''}}
\end{equation}

%%──────────────────────────────────────────────────────────────────
\subsection{Interpolation}
Interpolation is the art of reading ``between the lines'' of a discrete
dataset.
\begin{itemize}
    \item \textbf{Lagrange Interpolation}: Fits a single global polynomial of
    degree $N-1$ through exactly $N$ points. High-degree polynomials suffer
    from \textbf{Runge's Phenomenon}: they oscillate violently at the edges
    of the interval.
    \item \textbf{Cubic Splines}: Piecewise polynomials that enforce
    continuity of the function value, first derivative, and second derivative
    across each knot. The result is a stable tridiagonal system.
\end{itemize}

\begin{tcolorbox}[colback=green!5!white,colframe=green!50!black,
  title=The Tridiagonal System for Cubic Splines]
Let $h_i = x_{i+1} - x_i$ and $m_i = S''(x_i)$ be the unknown second
derivative at each knot. Enforcing continuity of $S$, $S'$, and $S''$ yields:
\begin{equation}
h_i\, m_{i-1} + 2(h_i + h_{i+1})\, m_i + h_{i+1}\, m_{i+1} =
6\left(\frac{y_{i+1}-y_i}{h_{i+1}} - \frac{y_i - y_{i-1}}{h_i}\right)
\end{equation}
With natural boundary conditions $m_0 = m_N = 0$, this gives an
$(N-1)\times(N-1)$ tridiagonal system solvable in $\bigO{N}$ — far cheaper
than $\bigO{N^3}$ Gaussian elimination.
\end{tcolorbox}

%%──────────────────────────────────────────────────────────────────
\subsection*{Recommended Reading}
\begin{itemize}
    \item \textit{Numerical Recipes}, W.H. Press et al.\ --- Ch.\ 1
    (Preliminaries on Errors), Ch.\ 5.3 (Polynomial Interpolation),
    Ch.\ 5.7 (Numerical Derivatives).
    \item \textit{Computational Physics}, Mark Newman --- Ch.\ 4 (Accuracy
    and Speed), Ch.\ 5 (Integrals and Derivatives).
\end{itemize}
